\usepackage{adjustbox}
\usepackage{algorithm}
\usepackage{algorithmic}
\usepackage{amsfonts}
\usepackage{amsmath}
\usepackage{array}
\usepackage{booktabs}
\usepackage{caption}
\usepackage{color}
\usepackage{enumitem}
\usepackage{fancyvrb}
\usepackage{float}
\usepackage[symbol]{footmisc}
\usepackage{inconsolata} % improve the aesthetics of text in
% the typewriter font
\usepackage{listings}
\usepackage{longtable}
\usepackage{makecell}
\usepackage{mathtools}
\usepackage{multirow}
% \usepackage{setspace}
\usepackage{subcaption}
\usepackage{supertabular}
\usepackage{tabto}
\usepackage{tablefootnote}
\usepackage{tabularx}
\usepackage{verbatim}
\usepackage{xfrac}
\usepackage{xltabular}

\renewcommand{\thefootnote}{\fnsymbol{footnote}}

\lstset{
xleftmargin=0pt, % remove indentation
basicstyle=\small\ttfamily,
columns=flexible,
breaklines=true,
aboveskip=0pt, % remove extra space above listings
belowskip=0pt, % remove extra space below listings
framexleftmargin=0pt, % remove indentation
}

\newcommand{\cy}[1]{}%{\textcolor{red}{[CY: #1]}}

%-------------------- environment definitions -----------------
\newenvironment{example}[0]{\begin{Example} \it}{\end{Example}}
\newenvironment{proof-of}[1]{{\em Proof of #1:}}{}

%----------------- command definitions ----------
\newcommand{\nnn}{{\tt \bf  not}}

% Any macro definitions you would like to include
% These are not defined in the style file, because they don't begin
% with \bmva, so they might conflict with the user's own macros.
% The \bmvaOneDot macro adds a full stop unless there is one in the
% text already.
\def\eg{\emph{e.g}\bmvaOneDot}
\def\Eg{\emph{E.g}\bmvaOneDot}
\def\etal{\emph{et al.}\bmvaOneDot}


\everypar{\looseness=-1}

\newcommand{\treelogo}{\raisebox{5pt}{\includegraphics[scale=0.050]{logos/tree.png}}}
\newcommand{\iiitdlogo}{\raisebox{5pt}{\includegraphics[scale=0.0113]{logos/iiitd-logo.png}}}
\newcommand{\ublogo}{\raisebox{5.2pt}
{\includegraphics[scale=0.085]{logos/ub-logo.png}}}
\newcommand{\fire}{\raisebox{5pt}{\includegraphics[scale=0.15]{logos/fire.png}}}
\newcommand{\adobelogo}{\raisebox{5pt}{\includegraphics[scale=0.032]{logos/adobe-logo.png}}}
\newcommand\coauth{$^\star$}
\newcommand\blfootnote[1]{%
  \begingroup
  \renewcommand\thefootnote{}\footnote{#1}%
  \addtocounter{footnote}{-1}%
  \endgroup
}

\usepackage[skins]{tcolorbox} % for little color boxes
\definecolor{ForestGreen}{RGB}{34,139,34}
\newcommand{\best}[1]{\colorbox{black!10}{#1}}
\definecolor{valbest}{HTML}{d9ead3}
\newcommand{\valbest}[1]{\colorbox{valbest}{#1}}
\definecolor{valgood}{HTML}{d0e0e3}
\newcommand{\valgood}[1]{\colorbox{valgood}{#1}}
\definecolor{valmid}{HTML}{fce5cd}
\newcommand{\valmid}[1]{\colorbox{valmid}{#1}}
\definecolor{valbad}{HTML}{ead1dc}
\newcommand{\valbad}[1]{\colorbox{valbad}{#1}}
\definecolor{themegreen}{HTML}{365956}
\definecolor{themepurple}{HTML}{3c1b48}
\definecolor{themered}{HTML}{b43748}

\newcolumntype{L}[1]{>{\raggedright\let\newline\\\arraybackslash\hspace{0pt}}m{#1}}
\newcolumntype{C}[1]{>{\centering\let\newline\\\arraybackslash\hspace{0pt}}m{#1}}
\newcolumntype{R}[1]{>{\raggedleft\let\newline\\\arraybackslash\hspace{0pt}}m{#1}}
\newcommand\Tstrut{\rule{0pt}{2.6ex}}         % = `top' strut
\newcommand\Bstrut{\rule[-1.3ex]{0pt}{0pt}}   % = `bottom' strut

\setlist[itemize]{noitemsep, topsep=0pt}
\setlist[enumerate]{noitemsep, topsep=0pt}

\newcommand{\citeauthorwithyear}[1]{\citeauthor{#1}~\shortcite{#1}}
\newcommand{\sota}{state-of-the-art}
\newcommand{\Sota}{State-of-the-art}

\newcommand{\parens}[1]{\left(#1\right)}
\newcommand{\braces}[1]{\left\{#1\right\}}
\newcommand{\bracks}[1]{\left[#1\right]}
\newcommand{\modulus}[1]{\left\vert#1\right\vert}
\newcommand{\norm}[1]{\left\Vert#1\right\Vert}
\newcommand{\angular}[1]{\langle#1\rangle}
\newcommand{\lmod}{\left|\!\left|}
\newcommand{\rmod}{\right|\!\right|}

\newcommand\numberthis{\addtocounter{equation}{1}\tag{\theequation}}

\DeclareMathOperator*{\minimize}{minimize}
\DeclareMathOperator*{\maximize}{maximize}

\newcommand{\samplei}{^{\parens{i}}}
\newcommand{\raisedth}[1]{#1^{\textrm{th}}}
% \linespread{0.94}

\newcommand{\accuracy}{$XX\%$}
